```latex
\documentclass[11pt]{article}

\usepackage{amsmath,amsthm,amssymb,mathtools}

\newtheorem{theorem}{Theorem}
\newtheorem{lemma}[theorem]{Lemma}

\begin{document}

\title{A cubic identity for three-term relations in hyperfields}
\author{}
\date{}
\maketitle

\begin{theorem}\label{thm:main}
Let $\mathbb H$ be a hyperfield with hyperaddition $\boxplus$, and let
$x,y,z\in \mathbb H$.  If
\[
0\in x\boxplus y\boxplus z,
\]
then
\[
0\in x^3\boxplus y^3\boxplus z^3\boxplus (-xyz)\boxplus (-xyz)\boxplus (-xyz).
\]
In particular, the answer to the question is \emph{yes}.
\end{theorem}

\medskip

We write $A\boxplus B:=\bigcup_{a\in A,\; b\in B}(a\boxplus b)$ for subsets
$A,B\subseteq \mathbb H$, and we use the same notation for iterated hypersums.
We also use the standard reversibility axiom for canonical hypergroups:
\[
u\in v\boxplus w \iff w\in u\boxplus (-v).
\]

We shall use the standard hyperproduct $\odot$ of polynomials over a hyperfield.
For polynomials
\[
f(T)=\sum_i a_iT^i,\qquad g(T)=\sum_j b_jT^j,
\]
the set $f\odot g$ consists of all polynomials $\sum_k c_kT^k$ such that
\[
c_k\in \boxplus_{i+j=k} a_ib_j
\qquad\text{for every }k.
\]
This hyperproduct is associative and commutative; see for example
\cite[Section~4]{Jun}.

\begin{lemma}[Fusion]\label{lem:fusion}
Let $A,B\subseteq \mathbb H$ and let $u\in \mathbb H$. If
\[
0\in A\boxplus u
\qquad\text{and}\qquad
0\in B\boxplus (-u),
\]
then
\[
0\in A\boxplus B.
\]
\end{lemma}

\begin{proof}
From $0\in A\boxplus u$, reversibility gives $-u\in A$.
From $0\in B\boxplus (-u)$, reversibility gives $u\in B$.
Hence
\[
0\in (-u)\boxplus u \subseteq A\boxplus B,
\]
as claimed.
\end{proof}

\begin{lemma}\label{lem:cubic-factor}
Assume that
\[
0\in x\boxplus y\boxplus z.
\]
Then there exists $a\in \mathbb H$ such that the cubic polynomial
\[
p(T):=T^3+aT-xyz
\]
lies in the triple hyperproduct
\[
(T-x)\odot (T-y)\odot (T-z).
\]
\end{lemma}

\begin{proof}
By reversibility, the hypothesis $0\in x\boxplus y\boxplus z$ implies
\[
x\in (-y)\boxplus (-z).
\]
Therefore
\[
q_x(T):=T^2+xT+yz \in (T-y)\odot (T-z),
\]
because the coefficient of $T$ in $(T-y)\odot (T-z)$ may be any element of
$(-y)\boxplus (-z)$.

Now compute $(T-x)\odot q_x(T)$. By definition of hyperproduct,
\[
(T-x)\odot q_x(T)
=
\Bigl\{
T^3+c_2T^2+c_1T-xyz
\;\Bigm|\;
c_2\in x\boxplus(-x),\;
c_1\in yz\boxplus(-x^2)
\Bigr\}.
\]
Since $0\in x\boxplus (-x)$, we may choose $c_2=0$. Choosing any
$a\in yz\boxplus (-x^2)$ then gives
\[
p(T)=T^3+aT-xyz \in (T-x)\odot q_x(T)
\subseteq (T-x)\odot (T-y)\odot (T-z).
\]
\end{proof}

\begin{lemma}\label{lem:root-relation}
Let $u,v,w\in \mathbb H^\times$ and $a\in \mathbb H$.
Assume that for some quadratic polynomial $q(T)$ one has
\[
T^3+aT-uvw \in (T-u)\odot q(T).
\]
Then
\[
0\in u^3\boxplus au\boxplus (-uvw).
\]
\end{lemma}

\begin{proof}
Write
\[
q(T)=dT^2+bT+c.
\]
Comparing coefficients in the relation
\[
T^3+aT-uvw \in (T-u)\odot (dT^2+bT+c),
\]
we obtain:
\begin{align*}
1 &= d,\\
0 &\in b\boxplus (-u),\\
a &\in c\boxplus (-ub),\\
-uvw &= (-u)c.
\end{align*}
Since $u\neq 0$, the multiplicative group law on $\mathbb H^\times$ lets us
cancel $-u$ in the last equality, hence $c=vw$.
Also, $0\in b\boxplus (-u)$ implies $b=u$ by uniqueness of additive inverses.
Therefore
\[
a\in vw\boxplus (-u^2).
\]
Multiplying by $u$ gives
\[
au\in uvw\boxplus (-u^3).
\]
By reversibility,
\[
au\in uvw\boxplus (-u^3)
\iff
0\in u^3\boxplus au\boxplus (-uvw),
\]
which is the desired conclusion.
\end{proof}

\begin{proof}[Proof of Theorem~\ref{thm:main}]
If one of $x,y,z$ is zero, the claim is immediate. Indeed, by symmetry,
assume $z=0$. Then $0\in x\boxplus y$, so $y=-x$, and hence
\[
0\in x^3\boxplus (-x)^3 = x^3\boxplus y^3.
\]
Since also $z^3=0$ and $xyz=0$, the required conclusion follows.

Thus we may assume from now on that
\[
x,y,z\in \mathbb H^\times.
\]

By Lemma~\ref{lem:cubic-factor}, there exists $a\in \mathbb H$ such that
\[
p(T):=T^3+aT-xyz \in (T-x)\odot (T-y)\odot (T-z).
\]
Because polynomial hypermultiplication is associative and commutative,
there exist quadratic polynomials $q_x(T),q_y(T),q_z(T)$ with
\[
p(T)\in (T-x)\odot q_x(T),\qquad
p(T)\in (T-y)\odot q_y(T),\qquad
p(T)\in (T-z)\odot q_z(T).
\]
Applying Lemma~\ref{lem:root-relation} three times (with
$(u,v,w)=(x,y,z)$, then $(y,x,z)$, then $(z,x,y)$) yields
\begin{align}
0&\in x^3\boxplus ax\boxplus (-xyz),\label{eq:rx}\\
0&\in y^3\boxplus ay\boxplus (-xyz),\label{eq:ry}\\
0&\in z^3\boxplus az\boxplus (-xyz).\label{eq:rz}
\end{align}

Next, multiply the original relation $0\in x\boxplus y\boxplus z$ by $-a$.
Using distributivity, we get
\[
0\in (-ax)\boxplus (-ay)\boxplus (-az).
\]
Now use Lemma~\ref{lem:fusion} repeatedly.

First fuse \eqref{eq:rx} with
$0\in (-ax)\boxplus (-ay)\boxplus (-az)$ along the pair $ax,-ax$. This gives
\[
0\in x^3\boxplus (-xyz)\boxplus (-ay)\boxplus (-az).
\]
Next fuse this relation with \eqref{eq:ry} along the pair $-ay,ay$. We obtain
\[
0\in x^3\boxplus y^3\boxplus (-xyz)\boxplus (-xyz)\boxplus (-az).
\]
Finally fuse this relation with \eqref{eq:rz} along the pair $-az,az$. This yields
\[
0\in x^3\boxplus y^3\boxplus z^3\boxplus (-xyz)\boxplus (-xyz)\boxplus (-xyz).
\]

This is exactly the desired conclusion.
\end{proof}

\begin{thebibliography}{9}

\bibitem{Jun}
J. Jun,
\emph{Algebraic geometry over hyperrings},
Adv. Math. \textbf{323} (2018), 142--192.

\end{thebibliography}

\end{document}
```